\begin{table}[htbp]
\centering
\caption{Preregistered Hypotheses Reconciliation and Final Research Claim Classifications}
\label{tab:table_T11_hypothesis_reconciliation}
\begin{tabular}{lllll}
\hline
Hypothesis & Original\_Prereg\_Target & Observed\_Result & Final\_Status & Paper\_Framing \\
\hline
H1: Hierarchical Pooling Advantage & Capped TTS reduction over independent SSM & Class-wide -7.8m (p=0.011), Broad -4.5m (p=0.011), Single +9.46m penalty & PARTIALLY\_SUPPORTED & Partial pooling trades broad-fleet acceleration for isolated-host penalty \\
H2: Two-Scale Cross-Level Inertia & Incremental benefit of hier\_ts over hier\_ssm & Mean diff 0.000 min, CI [0, 0] across 5 scenarios & NOT\_SUPPORTED & Negative finding: single-scale hierarchy captures all actionable dynamics \\
H3: Hysteresis Alert Reduction & >= 50\% non-event alert reduction with <= 3m delay & Max reduction 20.8\%, weak persistent delay +7.34m & NOT\_SUPPORTED & Negative finding: fixed hysteresis fails the 50\% threshold \\
H4: Recovery Dwell Acceleration & Faster return to normal after recovery & Recovery delay +4.81m; no acceleration observed & NOT\_SUPPORTED & Negative finding: recovery tracking incurs dwell cost rather than acceleration \\
H5: SSM Dominance over Classical Baselines & Universal superiority over CUSUM and GLR & Faster in 9/11 event scenarios; slower in slow drift (+19.6m) and step (+14.9m) & NO\_UNIVERSAL\_DOMINANCE & SSM is superior for abrupt/correlated faults; classical baselines win on slow drift \\
\hline
\end{tabular}
\end{table}
